%! Graphing Quadratic Functions & Transformations — Unit 2, Lesson 2.1 — Warm-Up
\documentclass[10pt]{article}
\usepackage{algebra2-article}
\usepackage{algebra2-boxes}
\newcommand{\TallMath}[1]{$\displaystyle #1\rule[-1.4em]{0pt}{3.2em}$}
\newcommand{\lgrid}[4]{%
  \draw[step=1, linegray!40, very thin] (#1,#3) grid (#2,#4);
  \draw[->, semithick, charcoal] (#1,0)--(#2,0) node[below right, font=\scriptsize]{$x$};
  \draw[->, semithick, charcoal] (0,#3)--(0,#4) node[left, font=\scriptsize]{$y$};}

\begin{document}

\pageheader{Unit 2, Lesson 2.1}{Warm-Up}

\vspace{0.10in}

\begin{notesbox}{Getting Ready: Skills We'll Use Today}
\small These skills power today's lesson on graphing parabolas. Work quickly.

\vspace{4pt}
\textbf{1. The parent $y=x^2$.} Complete the table of the five ``anchor'' points, then read the graph.
\begin{center}
\begin{minipage}{0.50\linewidth}\centering
\renewcommand{\arraystretch}{1.3}
\begin{tabular}{|c|c|c|c|c|c|}
\hline
$x$ & $-2$ & $-1$ & $0$ & $1$ & $2$ \\ \hline
$y=x^2$ & \blank{0.8cm} & \blank{0.8cm} & \blank{0.8cm} & \blank{0.8cm} & \blank{0.8cm} \\ \hline
\end{tabular}
\end{minipage}%
\begin{minipage}{0.40\linewidth}\centering
\begin{tikzpicture}[scale=0.42]
  \lgrid{-3}{3}{-1}{5}
  \begin{scope}
    \clip (-3,-1) rectangle (3,5);
    \draw[very thick, forest, domain=-2.2:2.2, samples=60, smooth] plot (\x,{(\x)*(\x)});
  \end{scope}
  \fill[charcoal] (0,0) circle (2.5pt);
\end{tikzpicture}
\end{minipage}
\end{center}
The vertex of $y=x^2$ is the point \blank{1.4cm}, and it opens \blank{1.2cm}.

\vspace{4pt}
\textbf{2. Move one point (Unit 1 review).} The point $(2,4)$ is on the parent $y=x^2$. Give the new
location of \emph{that} point after each change:
\[
  y=x^2+3:\ \blank{1.6cm} \qquad y=(x-1)^2:\ \blank{1.6cm} \qquad y=-x^2:\ \blank{1.6cm}
\]

\vspace{4pt}
\textbf{3. Expand a square.} Multiply out and simplify (write in the form $x^2+bx+c$).
\[
  (x-1)^2 - 4 = \blank{4.5cm}
\]
\end{notesbox}

\end{document}
